\chapter{Objectifs d'apprentissage}\label{V-objectifs-dapprentissage}

Cette section s'efforce d'expliciter et de hiérarchiser les acquis
d'apprentissages associés au chapitre. Ces objectifs sont organisés en
paliers~:

($\mathord{\boldsymbol{\circ}}$) Prérequis ($\mathord{\bullet}$)
Fondamental ($\mathord{\bullet}\mathord{\bullet}$) Standard
($\mathord{\bullet}\mathord{\bullet}\mathord{\bullet}$) Avancé
($\mathord{\bullet}\mathord{\bullet}\mathord{\bullet}\mathord{\bullet}$)
Expert

Sauf mention particulière, la connaissance des démonstrations du
document n'est pas exigible\footnote{l'étude des démonstrations du cours
  peut toutefois contribuer à votre apprentissage, au même titre que la
  résolution d'exercices.}

\runinhead{Intégrale de Monte
Carlo}\phantomsection\label{V-intuxe9grale-de-monte-carlo}

\begin{itemize}
\tightlist
\item
  $\mathord{\bullet}$ connaître le principe de l'intégration par la
  méthode Monte Carlo
\item
  $\mathord{\bullet}$ savoir que cette approche se justifie par la loi
  des grands nombres
\item
  $\mathord{\bullet}$ savoir que l'approximation fournie par le TCL
  offre un contrôle de l'erreur
\end{itemize}

\runinhead{Génération de nombres
pseudo-aléatoires}\phantomsection\label{V-guxe9nuxe9ration-de-nombres-pseudo-aluxe9atoires}

\begin{itemize}
\tightlist
\item
  $\mathord{\bullet}$ connaître le principe de la génération de nombres
  pseudo-aléatoires par la méthode des congruences
\end{itemize}

\runinhead{Méthodes de simulation de
v.a.}\phantomsection\label{V-muxe9thodes-de-simulation-de-v.a.}

\begin{itemize}
\tightlist
\item
  $\mathord{\bullet}$ connaître et savoir implémenter en python la
  méthode d'inversion
\item
  $\mathord{\bullet}$ connaître et savoir implémenter en python la
  méthode de rejet
\item
  $\mathord{\bullet}$ connaître la méthode de Box-Muller
\item
  $\mathord{\bullet}$ connaître et savoir implémenter la simulation de
  vecteurs gaussien par la méthode de Cholesky
\end{itemize}

\runinhead{Echantillonnage
d'importance}\phantomsection\label{V-echantillonnage-dimportance}

\begin{itemize}
\tightlist
\item
  $\mathord{\bullet}$ connaître la définition de la méthode
  d'échantillonnage d'importance
\item
  $\mathord{\bullet}$ savoir qu'un bon choix de densité instrumentale
  permet de minimiser la variance d'estimation
\end{itemize}

\chapter{Intégrale de
Monte-Carlo}\label{V-intuxe9grale-de-monte-carlo-1}

Les méthodes de Monte-Carlo ont été développées initialement par des
physiciens dans les années 1950 (notamment par les travaux de
Metropolis, Ulam, Hastings, Rosenbluth) pour calculer des intégrales
(déterministes) à partir de méthodes probabilistes numériquement assez
économiques. Le nom a été donné en référence au célèbre casino du fait
du caractère aléatoire de ces méthodes.

Les méthodes de simulation sont basées sur la production de nombres
aléatoires, distribués selon une certaine loi de probabilité. Dans de
nombreuses applications, pour une certaine fonction \(h\), on souhaite
calculer, pour une variable aléatoire \(X\) de loi \(\mathbb{P}_X\)
\[\mathcal{I}=\mathbb{E}\left(h(X)\right)=\int h(x)\, \mathbb{P}_X(dx).\]

En général, même si on sait évaluer \(h\) en tout point, on ne peut pas
calculer formellement l'intégrale \(\mathcal{I}\). Le calcul d'intégrale
par la méthode Monte-Carlo consiste dans sa version la plus simple à
générer un \emph{échantillon}
\((X_1,\ldots,X_n) \sim_{i.i.d.}\mathbb{P}_X\), et à estimer
\(\mathcal{I}\) par la moyenne empirique
\[M_n(h)=\frac{1}{n}\sum_{i=1}^{n}h(X_i),\] où i.i.d signifie
indépendant et identiquement distribué. En effet, d'après la
\href{Probabilité\%20III.pdf\%20\#lfgn}{loi forte des grands nombres},
si \(h(x)\) est \(\mathbb{P}_X\)-intégrable, on a l'assurance que quand
\(n \to +\infty\),
\[M_n(h) \rightarrow \int h(x)\mathbb{P}_X(dx) \text{ p.s.}\] Si de plus
\(h(X)^2\) est intégrable, la vitesse de convergence de \(M_n(h)\) peut
être évaluée, puisque la variance
\[\mathbb{V}(M_n(h)) = \frac{1}{n} \int \left(h(x)-\mathcal{I}\right)^2 \mathbb{P}_X(dx)\]
peut également être estimée à partir de l'échantillon
\((X_1,\ldots,X_n)\) par la quantité
\[\sigma_n^2=\frac{1}{n}\sum_{i=1}^{n}\left( h(X_i)-M_n(h)\right)^2.\]
Le \href{Probabilité\%20III.pdf\%20\#TCL}{théorème central limite}
assure alors que pour \(n\) grand,
\[\frac{M_n(h)-\mathcal{I}}{\sigma_n}\] suit approximativement une loi
\(\mathcal{N}(0,1)\)\footnote{ce résultat sera démontré dans le cours de
  science des données au second semestre.}. Cette propriété conduit à la
construction de tests de convergence et de bornes de confiance
asymptotiques pour \(M_n(h)\). Par exemple, on aura
\[\mathbb{P}\left(\mathcal{I} \in \left[M_n(h) - 1.96 \sigma_n,M_n(h) + 1.96 \sigma_n\right]\right) \approx 0.95,\]
où \(1.96 \approx \Phi^{-1}(0,975)\) avec \(\Phi\) la fonction de
répartition de la loi normale centrée réduite.

En outre, cette propriété indique que la vitesse de convergence de
\(M_n(h)\) est de l'ordre de \(\sqrt{n}\), et ce indépendamment de la
dimension du problème. Cela explique la supériorité de cette méthode par
rapport aux méthodes d'intégration numérique déterministes dont les
vitesses de convergence décroissent rapidement (exponentiellement) avec
la dimension du problème.

\runinhead{Exercice -- Calcul numérique d'une
probabilité (\(\mathord{\bullet}\mathord{\bullet}\))}\phantomsection\addcontentsline{toc}{subsubsection}{Calcul
numérique d'une probabilité}\label{V-mc1}
Calculer une estimation de la probabilité de l'événement
\(\sin(X) > 1/2\) pour \(X\) de loi \(\mathcal{N}(0, 1)\).

Indication~: si x est un vecteur,
\texttt{y=(np.sin(x)\textgreater{}1/2)} désigne le vecteur de booléens
dont les coordonnées valent \texttt{True} ou \texttt{False} selon que
les coordonnées correspondantes de \texttt{np.sin(x)} soient ou non
\textgreater{} 1/2. Par ailleurs, les vecteurs de booléens se
convertissent naturellement en vecteurs de 0 et de 1 si on leur applique
une fonction dont les arguments sont des nombres (ex~:
\texttt{np.mean}).\solref{V-answer-mc1}

\chapter{Génération de nombres
pseudo-aléatoires}\label{V-guxe9nuxe9ration-de-nombres-pseudo-aluxe9atoires-1}

Les ordinateurs sont des machines déterministes. Il peut sembler
paradoxal de leur demander de générer des nombres aléatoires\footnote{Von
  Neumann (1951) résume ce problème très clairement~: ``Any one who
  considers arithmetical methods of reproducing random digits is, of
  course, in a state of sin. As has been pointed out several times,
  there is no such thing as a random number --- there are only methods
  of producing random numbers, and a strict arithmetic procedure of
  course is not such a method.''}. En réalité, les algorithmes de
génération de nombres aléatoires vont produire des séquences de nombres
déterministes qui vont avoir l'aspect de l'aléatoire. On s'intéresse ici
à la génération de nombres uniformes sur \(]0,1[\) (on exclut les bornes
par commodité), puisque, comme on le verra dans la suite, il s'agit de
l'ingrédient de base de toutes les méthodes de simulation stochastique.

\addcontentsline{toc}{subsection}{Générateur de nombres uniformes
pseudo-aléatoires}
\begin{definition}[{Générateur de nombres uniformes
pseudo-aléatoires}]\phantomsection\label{V-guxe9nuxe9rateur-de-nombres-uniformes-pseudo-aluxe9atoires}

Un \emph{générateur de nombres uniformes pseudo-aléatoires} est un
algorithme qui étant donné une valeur initiale \(u_0\) et une
transformation \(T\) produit une séquence \(u_i=T^i(u_0)\),
\(i \in\mathbb{N}^\ast\), de valeurs dans \(]0,1[\).

Pour tout \(n\in\mathbb{N}^\ast\), les valeurs \((u_1,\ldots,u_n)\)
reproduisent le comportement d'une suite de variables aléatoires
\((V_1,\ldots,V_n)\) i.i.d de loi uniforme sur {]}0,1{[}, lorsqu'on les
compare au travers d'un ensemble de tests statistiques\footnote{Par
  exemple, la suite de tests
  \href{https://en.wikipedia.org/wiki/Diehard_tests}{Die Hard}, due à
  Marsaglia.}, vérifiant par exemple que la corrélation entre deux
nombres successifs est suffisamment faible.
\end{definition}

\addcontentsline{toc}{subsection}{la méthode des congruences}
\begin{exemple}[{la méthode des
congruences}]\phantomsection\label{V-la-muxe9thode-des-congruences}

Cet algorithme, dû à Lehmer (\protect\hyperlink{refV-Lehmer}{1951}), est
l'un des premiers à avoir été proposé et implémenté. Il repose sur 2
paramètres~:

\begin{itemize}
\tightlist
\item
  le multiplicateur \(a\in\mathbb{N}^\ast\),
\item
  le modulo \(m\in\mathbb{N}^\ast\).
\end{itemize}

Etant donné un entier \(k_0\) tel que \(1 \leq k_0 \leq m-1\), la
séquence de nombres est générée par la transformation suivante~: pour
tout \(n\in\mathbb{N}^\ast\) on définit l'entier
\(k_{n+1} = a\,k_n\,\mathrm{\,mod\,} m\) puis on pose
\[u_{n+1} = \frac{k_{n+1}}{m},\hspace{1em} u_0 = \frac{k_0}{m}.\]

On peut remarquer que l'algorithme va produire une séquence de valeurs
qui sera périodique, c'est-à-dire qu'après un certain nombre
d'itérations, la suite se répétera, et qu'il pourra fournir au plus
\(m-1\) valeurs différentes. Ce trait est commun à tous les générateurs
de nombre aléatoires et est lié aux limitations matérielles des
ordinateurs (on ne peut représenter qu'un nombre fini de nombres). Le
choix des valeurs de \(a\) et \(m\) est par conséquent crucial. Il
existe des critères qui permettent de s'assurer du bon comportement de
cette suite~:

\begin{itemize}
\tightlist
\item
  \(m\) est un nombre premier (le plus grand possible),
\item
  \(p = (m-1) /2\) est un nombre premier,
\item
  \(a^p = -1 \mathrm{\,mod\,} m\).
\end{itemize}

Ce critère assure une période pleine et donc que tous les nombres
\((1/m,\ldots,(m-1)/m)\) seront générés. Cependant, les nombres générés
ne sont pas indépendants, pas même non corrélés. On peut montrer que la
corrélation entre 2 nombres successifs vaut approximativement \(1/a\).
Il convient donc de choisir \(a\) suffisamment grand pour que celle-ci
devienne négligeable. Par exemple, en prenant \(a=1000\) et
\(m=2001179\), on obtient une période de \(2001178\) et une corrélation
de l'ordre de \(10^{-3}\).

Ce générateur de nombres uniformes pseudo-aléatoires, bien que
rudimentaire, a ouvert la voie à des générateurs plus sophistiqués,
toujours basés sur des opérations arithmétiques. Le plus couramment
utilisé à ce jour est l'algorithme de
\href{https://en.wikipedia.org/wiki/Mersenne_Twister}{Mersenne-Twister}.
Il s'agit du générateur par défaut de la plupart des logiciels tels que
Python, R, MATLAB, Julia, MS Excel,\ldots{} Il passe notamment avec
succès toute la batterie de tests Die Hard. En particulier, sa période
vaut \(2^{19937} - 1\).

Pour certains usages, cet algorithme n'est cependant pas recommandé du
fait de sa prédictibilité. C'est notamment un défaut rédhibitoire pour
les applications en cryptographie. Il existe des variantes mieux
adaptées à ce cas de figure. On notera enfin qu'il existe des
générateurs de nombres aléatoires basés sur des phénomènes physiques
comme un bruit électrique ou des phénomènes quantiques et donc
parfaitement imprévisibles.
\end{exemple}

\chapter{Méthodes de simulation de variables aléatoires
réelles}\label{V-muxe9thodes-de-simulation-de-variables-aluxe9atoires-ruxe9elles}

On a vu au chapitre II du cours de Probabilités que l'on pouvait
transformer des variables aléatoires réelles suivant certaines lois pour
en obtenir de nouvelles. Par exemple, si \(X_1,\dots,X_n\) sont
\(n\in\mathbb{N}^\ast\) variables gaussiennes centrées réduites
indépendantes, alors \(X_1^2+\dots,X_n^2\) suit une loi du \(\chi^2\) à
\(n\) degrés de liberté. Dans le même esprit, on va voir ici comment
simuler des v.a.r. de lois diverses à partir de la simulation de
variables uniformes sur \(]0,1[\). On introduit une notation qui sera
utile dans la suite~: pour spécifier que deux v.a.r. \(X\) et \(Y\) ont
même loi, on écrira \(X \overset{\mathcal{L}}{=} Y\).

\section{Méthode d'inversion}\label{V-muxe9thode-dinversion}

L'objectif de ce paragraphe est de définir quand et comment il est
possible de simuler une variable aléatoire réelle \(X\) de fonction de
répartition (f.d.r.) \(F_X\) en transformant la simulation d'une
variable aléatoire \(U\) de loi uniforme sur \(]0,1[\). En d'autres
termes, on cherche à déterminer les conditions sous lesquelles il est
possible d'identifier une fonction borélienne
\(\psi :\,]0,1[ \to \mathbb{R}\) telle que
\(X \overset{\mathcal{L}}{=} \psi(U)\).

Commençons par un cadre simple, où \(F_X\) est \textbf{bijective} d'un
intervalle non vide de \(\mathbb{R}\) sur \(]0,1[\).

\addcontentsline{toc}{subsection}{Cas où \(F_X\) est bijective}
\begin{proposition}[{Cas où \(F_X\) est
bijective}]\phantomsection\label{V-invbij}

Soient \(X\) une variable aléatoire réelle de fonction de répartition
\(F_X\) et \(U\) une variable uniforme sur \(]0,1[\). S'il existe un
intervalle non vide \(]a,b[ \subset \mathbb{R}\) tel que
\(F_X :\, ]a,b[ \to \,]0,1[\) est bijective, de bijection réciproque
\(F_X^{-1} :\, ]0,1[ \to ]a,b[\), alors
\(F_X^{-1}(U) \overset{\mathcal{L}}{=} X\) et
\(F_X(X) \overset{\mathcal{L}}{=} U\).
\end{proposition}

\begin{preuve}
Le premier résultat est immédiat~: pour tout \(x\in\mathbb{R}\), par
croissance de \(F_X\) et donc de \(F_X^{-1}\), on a
\[\mathbb{P}\left(F_X^{-1}(U)\leq x\right) = \mathbb{P}\left(U \leq F_X(x) \right) = F_X(x).\]
Concernant le second, notons \(G\) la fonction de répartition de la
variable aléatoire \(F_X(X)\). La croissance de \(F_X\) garantit que
\(\mathbb{P}\left(F_X(X) \in\, ]0,1[\right) = 1\). Ainsi, pour tout
\(x\in\mathbb{R}\) on a bien
\[G(x) = \left|\begin{array}{ll} 1 & \text{si } x \geq 1,\\ \mathbb{P}\left(F_X(X)\leq x\right) = \mathbb{P}\left( X \leq F_X^{-1}(x) \right) = x & \text{si } 0 < x < 1,\\ 0 & \text{sinon.}\end{array}\right.\]

\hfill$\blacksquare$
\end{preuve}

\runinhead{Exercice -- Exemples d'application
1 (\(\mathord{\bullet}\))}\phantomsection\addcontentsline{toc}{subsubsection}{Exemples
d'application 1}\label{V-exemples1}
Donner un algorithme de simulation d'une v.a.r. \(X\) suivant une loi

\begin{itemize}
\tightlist
\item
  uniforme sur un intervalle \(I \subset \mathbb{R}\),
\item
  Exponentielle de paramètre \(\lambda \in \mathbb{R}_+^\ast\),
\item
  de Cauchy, de densité
  \(x\in\mathbb{R}\mapsto \left(\pi\left(1+x^2\right)\right)^{-1}\),
\item
  de Laplace de paramètres \(\mu \in \mathbb{R}\) et
  \(s\in\mathbb{R}_+^\ast\), de densité
  \(x\in\mathbb{R}\mapsto \frac{1}{2s}\,\exp\left\{-\frac{|x-\mu|}{s}\right\}\),
\item
  Logistique de paramètres \(\mu \in \mathbb{R}\) et
  \(s\in\mathbb{R}_+^\ast\), de fonction de répartition
  \(x\in\mathbb{R}\mapsto \left(1 + \exp\left\{-\frac{x-\mu}{s} \right\}\right)^{-1}\)\,?
\end{itemize}

\solref{V-answer-exemples1}

Dans cette situation idéale, \(\psi = F_X^{-1}\) est une solution à
notre problème. Que se passe-t-il en revanche si \(F_X\) n'est pas
bijective\,?

\runinhead{Exercice -- Pile ou
face (\(\mathord{\bullet}\))}\phantomsection\addcontentsline{toc}{subsubsection}{Pile
ou face}\label{V-pf}
Proposer une méthode pour simuler un tir à pile ou face à partir de la
simulation d'une variable uniforme sur \(]0,1[\).\solref{V-answer-pf}

On a déjà vu au chapitre II que les fonctions de répartition de v.a.r.
possèdent un nombre au plus dénombrable de points de discontinuité. Sur
chaque intervalle où elles sont continues, on peut alors considérer
qu'elles sont bijectives, quitte à réduire les zones de palier à un
point. Cela permet de généraliser la notion de bijection réciproque pour
ces fonctions.

\addcontentsline{toc}{subsection}{Réciproque généralisée}
\begin{definition}[{Réciproque
généralisée}]\phantomsection\label{V-defrecgen}

Soit \(F\) une fonction de répartition. On définit sa \emph{réciproque
généralisée} (aussi appelée \emph{inverse généralisée} ou
\emph{pseudo-inverse}) comme la fonction
\[F^{-} : u \in\, ]0,1[\, \mapsto \inf\left\{ x \in \mathbb{R}: F(x) \geq u \right\} \in \mathbb{R}.\]
\end{definition}

\addcontentsline{toc}{subsection}{Conséquences}
\begin{remarque}[{Conséquences}]\phantomsection\label{V-consuxe9quences}

\begin{itemize}
\tightlist
\item
  Cette fonction est bien définie sur tout \(]0,1[\), car quel que soit
  \(u\) dans cet intervalle, l'ensemble
  \(\left\{ x \in \mathbb{R}: F(x) \geq u \right\}\) est non vide et
  minoré. S'il était vide ou non minoré pour un certain
  \(u_0\in\,]0,1[\), pour tout \(x \in \mathbb{R}\) on aurait dans le
  premier cas \(F(x) < u_0 < 0\) et dans le second
  \(F(x) \geq u_0 > 1\). L'une comme l'autre de ces inégalités est
  impossible pour une fonction de répartition.
\item
  La réciproque généralisée de la f.d.r. \(F_X\) d'une v.a.r. \(X\) est
  aussi appelée \emph{fonction quantile}. On pourra notamment remarquer
  que \(F_X^{-}\left(\frac{1}{2}\right)\) est la médiane de \(X\).
\item
  Lorsque \(F\) réalise une bijection d'un intervalle non vide
  \(I\subset \mathbb{R}\) sur \(]0,1[\), sa réciproque généralisée
  coïncide avec sa bijection réciproque.
\end{itemize}

On a alors le résultat suivant, qui indique que \(\psi = F_X^-\) est une
solution universelle à notre problème. La preuve détaillée est donnée en
Annexe.
\end{remarque}

\addcontentsline{toc}{subsection}{Méthode d'inversion}
\begin{theoreme}[{Méthode d'inversion}]\phantomsection\label{V-invgen}

Soient \(U\) une variable uniforme sur \(]0,1[\) ainsi que \(X\) une
variable aléatoire réelle de fonction de répartition \(F_X\) et de
réciproque généralisée \(F_X^-\). Alors
\(F_X^-(U) \overset{\mathcal{L}}{=} X\).
\end{theoreme}

\runinhead{Exercice -- Exemples d'application
2 (\(\mathord{\bullet}\))}\phantomsection\addcontentsline{toc}{subsubsection}{Exemples
d'application 2}\label{V-exemples2}
Donner un algorithme de simulation d'une v.a.r. \(X\) suivant une loi

\begin{itemize}
\tightlist
\item
  Binomiale de paramètres \(n\in\mathbb{N}^\ast\) et \(p \in\, ]0,1[\),
\item
  Uniforme sur l'union de deux segments non vides et disjoints
  \([a,b], [c,d]\subset\mathbb{R}\), de densité
  \(x\in\mathbb{R}\mapsto (b-a + d-c)^{-1}\,1_{[a,b]\cup[c,d]}(x)\)\,?
\end{itemize}

\solref{V-answer-exemples2}

\subsection{Limitations}\label{V-limitations}

La méthode d'inversion peut sembler universelle pour simuler une v.a.r.
\(X\) à partir de \(U \sim \mathcal{U}_{]0,1[}\). Cependant, elle
nécessite en pratique de disposer d'une expression analytique de \(F_X\)
pour en déduire sa réciproque généralisée. Ce n'est typiquement pas le
cas de nombreuses lois usuelles comme la loi Normale. On va donc
déterminer d'autres procédures pour simuler des variables suivant de
telles lois.

\section{Méthode de rejet}\label{V-muxe9thode-de-rejet}

La méthode de rejet est une alternative populaire à la méthode
d'inversion, lorsque cette dernière ne peut être utilisée directement et
que \textbf{la loi cible possède une densité}. On la doit à Von Neumann
(\protect\hyperlink{refV-vonNeumann}{1951}). Pour en comprendre le
fondement, il faut d'abord introduire une généralisation naturelle de la
loi uniforme dans \(\mathbb{R}\) à \(\mathbb{R}^d\)
(\(d\in\mathbb{N}^\ast\)). On notera \(\lambda\) la mesure de Lebesgue
sur \(\mathbb{R}^d\).

\addcontentsline{toc}{subsection}{Loi uniforme sur un borélien}
\begin{definition}[{Loi uniforme sur un
borélien}]\phantomsection\label{V-loi-uniforme-sur-un-boruxe9lien}

La loi uniforme sur un borélien \(A\subset\mathbb{R}^d\) de volume
\(\lambda(A) > 0\) est une loi de probabilité admettant pour densité
\[f : x\in\mathbb{R}^d \mapsto \dfrac{1_A(x)}{\lambda(A)}.\]
\end{definition}

\runinhead{Exercice -- Loi uniforme sur un
pavé (\(\mathord{\bullet}\))}\phantomsection\addcontentsline{toc}{subsubsection}{Loi
uniforme sur un pavé}\label{V-unipave}
Comment simuler un vecteur aléatoire \((U_1,\dots,U_d)\) de loi uniforme
sur un pavé non vide
\([a_1,b_1]\times\dots\times[a_d,b_d] \subset \mathbb{R}^d\)\,?\solref{V-answer-unipave}

Une propriété fondamentale de cette loi est qu'elle reste stable par
conditionnement, dans le sens suivant.

\addcontentsline{toc}{subsection}{Stabilité par conditionnement}
\begin{proposition}[{Stabilité par
conditionnement}]\phantomsection\label{V-stabunif}

Soit \(U\) un vecteur aléatoire de loi uniforme sur un borélien
\(A\subset\mathbb{R}^d\) de volume \(\lambda(A) > 0\). Alors pour tout
borélien \(B \subset A\) de volume \(\lambda(B) > 0\), la loi
conditionnelle \(\mathbb{P}_{U\mid U\in B}\) de \(U\) sachant que
\(U \in B\) est uniforme sur \(B\).
\end{proposition}

\runinhead{Exercice -- Démonstration (\(\mathord{\bullet}\))}\phantomsection\addcontentsline{toc}{subsubsection}{Démonstration}\label{V-unicond}
\solref{V-answer-unicond}

\runinhead{Exercice -- Simulation (\(\mathord{\bullet}\))}\phantomsection\addcontentsline{toc}{subsubsection}{Simulation}\label{V-simunibor}
Déduire de la \protect\hyperref[V-stabunif]{proposition
précédente (p.~\pageref*{V-stabunif})} une méthode pour simuler un
vecteur aléatoire de loi uniforme sur un borélien
\(B \subset \mathbb{R}^d\) \textbf{borné} et de volume
\(\lambda(B) > 0\).\solref{V-answer-simunibor}

Il est aussi possible de simuler un vecteur aléatoire de loi uniforme
sur certains boréliens non vides et non bornés, comme l'établit la
proposition ci-dessous.

\addcontentsline{toc}{subsection}{Loi uniforme sur le sous-graphe d'une
densité}
\begin{proposition}[{Loi uniforme sur le sous-graphe d'une
densité}]\phantomsection\label{V-simunifdens}

Soient une densité \(f : \mathbb{R}\to \mathbb{R}\), une v.a.r. \(X\) et
une variable \(U\) uniforme sur \(]0,1[\), indépendante de \(X\). On
note
\(A_f := \left\{ (x,y) \in \mathbb{R}\times \mathbb{R}_+ : f(x) \geq y \right\}\)
le domaine limité par le graphe de \(f\) et l'axe des abscisses (souvent
appelé sous-graphe de \(f\)). Si \(X\) est de densité \(f\), alors le
couple \((X,Uf(X))\) suit une loi uniforme sur \(A_f\). Réciproquement, si \((X,Y)\) suit une loi uniforme sur \(A_f\), alors \(X\) est de densité \(f\) et \(Y\) conditionnellement à \(X\) est uniforme sur \(]0,f(X)[\).
\end{proposition}

\begin{preuve}
Commençons par remarquer que \(\lambda(A_f) = 1\). Quel que soit
\((z,v)\in\mathbb{R}^2\), par indépendance de \(X\) et \(U\) et par
Fubini on a \begin{align*}
\mathbb{P}\left( X \leq z, Uf(X) \leq v \right) &= \int_\mathbb{R}\int_\mathbb{R}1_{]-\infty,z]}(x)\,1_{]-\infty,v]}(uf(x))\,1_{]0,1[}(u)\,f(x)\,du\,dx\\
&= \int_{-\infty}^z \left(\int_\mathbb{R}1_{]-\infty,v]\cap ]0,f(x)[}(uf(x))\,f(x)\,du\right)\,dx\\
&= \int_{-\infty}^z \int_{-\infty}^v 1_{]0,f(x)[}(u)\,du\,dx\\
& = \int_{-\infty}^z \int_{-\infty}^v 1_{A_f}(x,u)\,du\,dx.
\end{align*} Ainsi, \((X,Uf(X))\) admet pour densité \(1_{A_f}\), qui
correspond bien à celle d'une loi uniforme sur
\(A_f\).

Réciproquement, supposons que \((X,Y)\) suit une loi uniforme sur
\(A_f\). Alors pour tout \(x\in\mathbb{R}\)
\[\mathbb{P}\left(X \leq x\right) = \int_{-\infty}^x \int_\mathbb{R} 1_{A_f}(u,v)\,du\,dv = \int_{-\infty}^x f(u)\,du,\] ce qui montre que \(X\) est de densité \(f\). De plus, pour tout \(x\in\mathbb{R}\) et \(y\in\mathbb{R}_+\),
\begin{align*}
\mathbb{P}\left(Y \leq y, X \leq x\right) 
&= \int_{-\infty}^x \int_\mathbb{R} 1_{A_f}(u,v)\,1_{]-\infty,y]}(v)\,dv\,du\\
&= \int_{-\infty}^x \int_0^{\min(f(u),y)} dv\,du\\
&= \int_{-\infty}^x \min(f(u),y)\,du.
\end{align*}
ce qui montre que \(Y\) conditionnellement à \(X\) est uniforme sur \(]0,f(X)[\).

\hfill$\blacksquare$
\end{preuve}

Pour simuler un vecteur uniforme \((X,Y)\) sur un ensemble \(A_f\) tel
que défini à la proposition précédente, il suffit donc de simuler une
v.a.r. \(X\) de densité \(f\), puis une variable \(U\) uniforme sur
\(]0,1[\), et de poser \(Y = f(X)U\).

En combinant les résultats précédents, on obtient la méthode de rejet,
illustrée sur la figure ci-dessous, où
\(A_Y := \left\{ (x,y) \in \mathbb{R}\times \mathbb{R}_+ : y \leq f_Y(x) \right\}\).

\subsection{Méthode de rejet}\label{V-muxe9thode-de-rejet-1}

On souhaite simuler une variable aléatoire réelle \(X\) de densité
\(f_X\). Supposons que l'on sait simuler une variable \(U\) uniforme sur
\(]0,1[\), ainsi qu'une v.a.r. \(Y\) (par exemple avec la méthode
d'inversion) de densité \(f_Y\) telle qu'il existe un réel \(a > 0\)
pour lequel on a \(\forall x \in \mathbb{R}\)~:
\(f_X(x) \leq a\,f_Y(x)\). Il suffit alors de suivre l'algorithme
suivant~:

\begin{enumerate}
\def\labelenumi{\arabic{enumi}.}
\item
  simuler \(Y\) et \(U\),
\item
  si \(aUf_Y(Y) > f_X(Y)\), recommencer à l'étape 1.
\item
  poser \(X = Y\).
\end{enumerate}

\begin{figure}
\centering
\includegraphics{MetRej.tex.pdf}
\caption{Méthode de rejet}\label{V-fig-MetRej.tex}
\end{figure}

\subsection{Limitations}\label{V-limitations-1}

La méthode de rejet a l'avantage de permettre de simuler des variables
aléatoires à densité lorsque la méthode d'inversion est inapplicable. Néanmoins, pour pouvoir l'appliquer il faut indiquer une densité instrumentale qui,
multipliée par un réel positif, majore la densité cible, et que l'on
sait simuler. Le taux de rejet, c'est-à-dire la probabilité de
l'événement \(\{aUf_Y(Y) > f_X(Y)\}\), peut parfois être élevé,
notamment lorsque la dimension de \(X\) est grande, ce qui limite
l'efficacité de la méthode.

\runinhead{Exercice -- Taux de
rejet (\(\mathord{\bullet}\mathord{\bullet}\))}\phantomsection\addcontentsline{toc}{subsubsection}{Taux
de rejet}\label{V-tauxrej}
Calculer le taux de rejet de la méthode proposée
ci-dessus.\solref{V-answer-tauxrej}

\section{Simulation de variables aléatoires gaussiennes~:
Box-Muller}\label{V-simulation-de-variables-aluxe9atoires-gaussiennes-box-muller}

On a vu que la méthode d'inversion est inappropriée pour simuler une
variable gaussienne, puisqu'elle requiert l'expression analytique de la
fonction de répartition cible. La méthode de rejet est quant à elle sous-optimale,dans le sens où toutes les variables uniformes générées ne sont pas directement utilisées.
La loi normale étant fondamentale en probabilité, il est nécessaire de disposer d'une méthode de simulation exacte et
efficace.

George E. P. Box et Mervin E. Muller ont proposé en 1958 une telle
méthode (G.E.P. Box and M.E. Muller
(\protect\hyperlink{refV-BoxMuller}{1958})). Elle exploite la propriété
d'invariance par rotation de la densité d'un couple de variables
gaussiennes indépendantes centrées réduites.

\begin{proposition}\phantomsection\label{V-boxmuller}
Soient \(U\) et \(V\) deux variables indépendantes, de loi uniforme sur
\(]0,1[\). Alors les variables aléatoires
\(X = \sqrt{-2\ln(U)}\cos\left(2\pi V\right)\) et
\(Y = \sqrt{-2\ln(U)}\sin\left(2\pi V\right)\) sont indépendantes et
suivent toutes deux une loi normale centrée réduite.
\end{proposition}

\begin{preuve}
On considère \((\widetilde{X},\widetilde{Y})\) un vecteur aléatoire dont
les deux composantes sont indépendantes, de loi normale centrée réduite.
On note ses coordonnées polaires aléatoires \(\widetilde{R}\) et
\(\widetilde{\Theta}\). On a vu dans le cours de Probabilités II que
dans ce cas, \(\widetilde{R}\) et \(\widetilde{\Theta}\) sont
indépendantes, la première de densité
\(f_{\widetilde{R}} : r \in \mathbb{R}\mapsto r\,e^{-\frac{r^2}{2}} 1_{\mathbb{R}_+^\ast}(r)\)
et la seconde de loi uniforme sur \(]0,2\pi]\).

\(X\) et \(Y\) sont les coordonnées cartésiennes obtenues à partir d'un
rayon et d'un angle~: en posant \(R = \sqrt{-2\ln(U)}\) et
\(\Theta = 2\pi V\) on obtient \(X =R\cos(\Theta)\) et
\(Y = R\sin(\Theta)\). Par indépendance de \(U\) et \(V\), \(R\) et
\(\Theta\) sont indépendantes. Pour que le vecteur \((X,Y)\) ait la même
distribution que
\((\widetilde{X},\widetilde{Y}) = \left(\widetilde{R}\cos(\widetilde{\Theta}),\widetilde{R}\sin(\widetilde{\Theta})\right)\),
il suffit de montrer que \(R \overset{\mathcal{L}}{=} \widetilde{R}\) et
\(\Theta \overset{\mathcal{L}}{=} \widetilde{\Theta}\).

\begin{itemize}
\item
  Commençons par étudier la loi de \(R\), de fonction de répartition
  \(F_R\). On remarque que la fonction
  \(u\in\, ]0,1[ \mapsto \sqrt{-2\ln(u)} \in \mathbb{R}_+^\ast\) est
  bijective, strictement décroissante. Ainsi, pour tout
  \(r\in\mathbb{R}_-\) on a \(\mathbb{P}\left(R \leq r \right) = 0\) et
  pour tout \(r \in \mathbb{R}_+^\ast\) on a
  \[\mathbb{P}(R \leq r) = \mathbb{P}\left(\sqrt{-2\ln(U)} \leq r \right) = \mathbb{P}\left(U \geq e^{-\frac{r^2}{2}} \right) = 1 - e^{-\frac{r^2}{2}}.\]
  qui correspond à la fonction de répartition associée à la densité
  \(f_{\widetilde{R}}\).
  \[\int_{-\infty}^r f_{\widetilde{R}}(x)\,dx = \left|\begin{array}{ll}\displaystyle \int_0^r x\,e^{-\frac{x^2}{2}}\,dx = \left[-e^{-\frac{x^2}{2}} \right]_0^r = 1 - e^{-\frac{r^2}{2}}  & \text{si } r>0,\\[1em] 0 & \text{sinon.} \end{array}\right.\]
\item
  En ce qui concerne la loi de \(\Theta\), de fonction de répartition
  \(F_\Theta\), puisque la fonction
  \(v \in\, ]0,1[\, \mapsto 2\pi v \in\, ]0,2\pi[\) est bijective
  strictement croissante, on a que pour tout \(\theta \in \mathbb{R}\)
  \[F_\Theta(\theta) = \left|\begin{array}{ll} 1 & \text{si } \theta \geq 2\pi,\\ \mathbb{P}\left(V\leq \frac{\theta}{2\pi} \right) = \dfrac{\theta}{2\pi} & \text{si } \theta \in\, ]0,2\pi[,\\ 0 & \text{si } \theta \leq 0,\end{array}\right.\]
  qui est bien la fonction de répartition d'une loi uniforme sur
  \(]0,2\pi[\).
\end{itemize}

\hfill$\blacksquare$
\end{preuve}

Cette méthode permet de simuler directement deux variables gaussiennes
centrées réduites indépendantes à partir de deux variables uniformes
indépendantes. Pour simuler une variable gaussienne d'espérance
\(m \in \mathbb{R}\) et de variance \(\sigma^2 \in \mathbb{R}_+^\ast\)
quelconques, on se rappelera que si \(X\) suit une loi normale centrée
réduite, alors \(\sigma X + m\) suit une loi normale d'espérance \(m\)
et de variance \(\sigma^2\). Dans les logiciels de calcul scientifique (python, R, Julia, etc.), la fonction de simulation d'une variable gaussienne repose cependant sur d'autres algorithmes, plus rapides, comme l'algorithme de la \emph{ziggurat} (G. Marsaglia and W. W. Tsang
(\protect\hyperlink{refV-Ziggurat}{2000})) : on recouvre la densité par
des rectangles empilés de même aire, on tire un point uniformément dans
un rectangle, et on ne recourt à un calcul coûteux (exponentielle, rejet) que dans de rares cas.

\chapter{Simulation d'un vecteur gaussien à
densité}\label{V-simulation-dun-vecteur-gaussien-uxe0-densituxe9}

On souhaite simuler un vecteur gaussien \(X = (X_1,\ldots,X_d)\) à
valeurs dans \(\mathbb{R}^d\) d'espérance \(m\) et de matrice de
covariance \(C\) définie positive (et donc inversible) données.

Puisque la matrice \(C\) est définie positive, elle admet une racine
carrée, c'est-à-dire qu'il existe une matrice \(N\) telle que
\(C = N\,N^t\). En effet, on peut par exemple décomposer \(C\) de la
manière suivante~: \[C = V\,D\,V^t,\] où \(V\) est une matrice
orthogonale et \(D\) est la matrice diagonale dont les termes diagonaux
sont les valeurs propres (toutes strictement positives) de \(C\). Il
suffit alors de prendre \(N = V\,D^{1/2}\), où \(D^{1/2}\) est la
matrice diagonale dont les termes diagonaux sont les racines carrées des
valeurs propres.

En pratique, il est coûteux numériquement d'effectuer le calcul des
valeurs propres et des vecteurs propres de \(C\). On va plutôt calculer
sa
\href{https://fr.wikipedia.org/wiki/Factorisation_de_Cholesky}{\emph{décomposition}
ou \emph{factorisation de Cholesky}} qui permet d'écrire \[C = L\,L^t\]
avec \(L\) une matrice triangulaire inférieure.\footnote{Cette
  décomposition est très utile dans la résolution de systèmes linéaires
  de la forme \(A\,x = b\), où \(b\) est connu, \(x\) inconnu et \(A\)
  est définie positive. Cela revient à résoudre \(L\,L^t\,x = b\). On
  pose alors \(y = L^t\,x\) et on résout d'abord \(Ly=b\), ce qui est
  très rapide puisque \(L\) est triangulaire inférieure (on commence par
  \(y_1 = b_1/L_{11}\), puis \(y_2 = (b_2 - L_{21}y_1)/L_{22}\), etc. en
  descendant). On résout ensuite \(L^t\,x = y\), ce qui est aussi très
  rapide pour la même raison (on commence par \(x_n = y_n/L_{nn}\) puis
  on remonte).}

Soit maintenant un autre vecteur gaussien \(Y = (Y_1,\ldots,Y_d)\) à
valeurs dans \(\mathbb{R}^d\) et de matrice de covariance identité,
notée \(I_d\). Autrement dit, les \(Y_i\) sont des variables aléatoires
gaussiennes centrées, réduites et indépendantes.

Alors, le vecteur \(Z = m + L\,Y\) est gaussien, d'espérance \(m\) et de
matrice de covariance \(C\). En effet, \(Z\) est gaussien comme
combinaison linéaire de variables aléatoires gaussiennes,
\(\mathbb{E}(Z) = \mathbb{E}(m + L\,Y) = m\) et
\(\mathbb{V}(Z) = \mathbb{E}((L\,Y)^2) = L I_d L^t = C\).

\chapter{Echantillonnage
d'importance}\label{V-echantillonnage-dimportance-1}

On introduit dans cette section la méthode d'échantillonnage
d'importance (importance sampling en anglais), que l'on appelle aussi,
de manière plus intuitive, échantillonnage préférentiel, pour les lois à
densité. Pour ce faire, nous allons commencer par un exemple qui montre
qu'il peut être plus efficace de simuler des valeurs selon une loi
différente de celle d'intérêt, autrement dit de modifier la
représentation de l'intégrale \(\mathcal{I}\) sous la forme d'une
espérance calculée selon une autre densité.

\addcontentsline{toc}{subsection}{Loi de Cauchy conditionnelle}
\begin{exemple}[{Loi de Cauchy
conditionnelle}]\phantomsection\label{V-loi-de-cauchy-conditionnelle}

Supposons que l'on s'intéresse à calculer la probabilité \(p\) qu'une
variable \(X\) de loi de Cauchy standard soit plus grande que 2 (on peut
le calculer directement et \(p=0.15\))
\[p = \int_2^{+\infty} \frac{1}{\pi(1+x^2)}dx.\] Si on estime \(p\)
directement à partir d'un échantillon \((X_1,\ldots,X_n)\) simulé selon
la loi de Cauchy standard, soit
\[\widehat{p}_1 = \frac{1}{n}\sum_{i=1}^n 1_{X_i > 2},\] la
\emph{variance de l'estimateur}
\(\mathbb{V}(\widehat{p}_1) = p(1-p)/n = 0.127/n\), puisque
\(\widehat{p}_1\) suit une loi binomiale de paramètre \((n,p)\). On peut
réduire cette variance (et donc améliorer la qualité de l'estimateur) en
tirant parti de la symétrie de la densité de la loi de Cauchy, en
formant un second estimateur
\[\widehat{p}_2 = \frac{1}{n}\sum_{i=1}^n 1_{|X_i| > 2},\] dont la
variance vaut \(\mathbb{V}(\widehat{p}_2) = p(1-p/2)/2n = 0.052/n\). La
relative inefficacité de ces méthodes est due au fait que la majeure
partie des valeurs simulées seront en dehors de la zone d'intérêt
\(]2,+\infty[\). En passant par le complémentaire, on peut réécrire
\(p\) comme \[p = \frac{1}{2} - \int_0^2 \frac{1}{\pi(1+x^2)}dx,\] dont
le second terme peut être vu comme l'espérance de
\(h(U) = \frac{2}{\pi(1+U^2)}\) avec \(U\sim \mathcal{U}_{[0,2]}\).
Tirant un échantillon \((U_1,\ldots,U_n)\) i.i.d. de loi uniforme sur
\([0,2]\), on obtient un troisième estimateur~:
\[\widehat{p}_3 = \frac{1}{2} - \frac{1}{n}\sum_{i=1}^n \frac{2}{\pi(1+U_i^2)},\]
dont la variance vaut
\(\mathbb{V}(\widehat{p}_3) = (\mathbb{E}(h(X)^2) - \mathbb{E}(h(U))^2)/n = 0.0285/n\)
(par intégration par parties). Enfin, on peut encore réécrire (voir
Ripley (\protect\hyperlink{refV-ripley}{1987}))
\[p = \int_0^{1/2}\frac{y^{-2}}{\pi(1+y^{-2})}dy,\] qui peut être vue
comme \(\mathbb{E}\left(\frac{V^{-2}}{2\pi(1+V^{-2})}\right)\) avec
\(V\sim \mathcal{U}_{]0,1/2[}\). L'estimateur formé à partir de cette
représentation et d'un échantillon \((V_1,\ldots,V_n)\) i.i.d. de loi
uniforme sur \([0,1/2]\) a une variance de \(0.95\, 10^{-4}/n\). Il est
donc bien plus efficace que \(\widehat{p}_1\) puisqu'il nécessite
environ \(\sqrt{10^3}\approx 32\) fois moins de simulations pour
atteindre la même précision.

On a ainsi vu sur ce cas particulier que l'estimation d'une intégrale de
la forme
\[\mathcal{I}=\mathbb{E}\left(h(X)\right)=\int_{\mathbb{R}^d} h(x) f(x) dx,\]
peut s'écrire de différentes manières, en faisant varier \(h\) et \(f\).
Par conséquent, un estimateur ``optimal'' devrait tenir compte de
l'ensemble de ces possibilités. C'est justement l'idée développée dans
la méthode d'échantillonnage d'importance dont le principe est décrit
dans la définition suivante.
\end{exemple}

\addcontentsline{toc}{subsection}{Echantillonnage d'importance}
\begin{definition}[{Echantillonnage
d'importance}]\phantomsection\label{V-is}

La méthode d'\emph{échantillonnage d'importance} est une évaluation de
\(\mathcal{I}\) basée sur la simulation d'un échantillon
\(X_1, \ldots,X_n\) de loi de densité \(g\) et approximant~:
\[\mathbb{E}_{f}\left(h(X)\right)\approx \frac{1}{n}\sum_{i=1}^{n}\frac{f(X_i)}{g(X_i)}h(X_i),\]
où la notation \(\mathbb{E}_f\) signifie que l'espérance est calculée
avec \(X\sim f\). On appelle souvent les ratios
\(\frac{f(X_i)}{g(X_i)}\) les \emph{poids d'importance} que l'on note
\(w_i\).
\end{definition}

Cette méthode est basée sur la représentation suivante de
\(\mathcal{I}\)~:

\[\mathbb{E}_{f}\left[h(X)\right]=\int h(x)\frac{f(x)}{g(x)}g(x)\,dx\]

que l'on appelle l'\emph{identité fondamentale de l'échantillonnage
d'importance} et l'estimateur converge du fait de la
\protect\hyperref[II-lfgn]{loi forte des grands nombres}.

Cette identité indique qu'une intégrale du type \(\mathcal{I}\) n'est
pas intrinsèquement associée à une loi donnée. L'intérêt de
l'échantillonnage d'importance repose sur le fait qu'il n'y a aucune
restriction sur le choix de la densité \(g\), dite \emph{instrumentale},
que l'on peut donc choisir parmi les densités des lois que l'on sait
simuler aisément. Il y a bien évidemment des choix qui sont meilleurs
que d'autres. Remarquons tout d'abord que bien que l'estimateur proposé
dans la \protect\hyperref[V-is]{définition
ci-dessus (p.~\pageref*{V-is})} converge presque sûrement, sa variance
est finie si
\[\mathbb{E}_g\left(h^2(X)\frac{f^2(X)}{g^2(X)}\right) = \mathbb{E}_f\left(h^2(X)\frac{f(X)}{g(X)}\right) = \int h^2(x)\frac{f^2(x)}{g(x)}\,dx < \infty.\]
On préconise alors l'usage de densités instrumentales \(g\) dont la
queue de distribution est plus épaisse que celle de \(f\) pour éviter
que cette variance puisse être infinie (on notera que cela dépend aussi
de la fonction \(h\) à intégrer). En pratique, on utilise généralement
l'estimateur suivant, de variance finie et qui donne des résultats plus
stables numériquement que celui de la définition~:
\[\frac{\sum_{i=1}^{n}w_ih(X_i)}{\sum_{i=1}^n w_i}\] où on a remplacé
\(n\) par la somme des poids d'importance. Puisque
\(\frac{1}{n}\sum_{i=1^n}w_i = \frac{1}{n}\sum_{i=1}^n\frac{f(x_i)}{g(x_i)}\)
tend vers 1 quand \(n\to\infty\), cet estimateur converge presque
sûrement vers \(\mathbb{E}_{f}\left(h(X)\right)\) par la loi forte des
grands nombres (voir C. P. Robert and G. Casella
(\protect\hyperlink{refV-roca}{2004})).

Parmi les densités \(g\) qui fournissent des estimateurs de variance
finie, il est possible d'exhiber la densité optimale (au sens de la
variance de l'estimateur associé) pour une fonction \(h\) et une densité
\(f\) données.

\addcontentsline{toc}{subsection}{Densité optimale}
\begin{theoreme}[{Densité
optimale}]\phantomsection\label{V-densituxe9-optimale}

Le choix de \(g\) qui minimise la variance de l'estimateur donné dans la
\protect\hyperref[V-is]{définition ci-dessus (p.~\pageref*{V-is})} est
\[g^\ast(x) = \frac{|h(x)|f(x)}{\int |h(x)|f(x) dx}.\]
\end{theoreme}

\begin{preuve}
Notons d'abord que
\[\mathbb{V}_g\left(\frac{h(X)f(X)}{g(X)}\right) = \mathbb{E}_g\left(h^2(X)\frac{f^2(X)}{g^2(X)}\right) - \mathbb{E}_g\left(h(X)\frac{f(X)}{g(X)}\right)^2\]
et le second terme ne dépend pas de \(g\). On minimise donc le premier
terme. D'après \protect\hyperref[II-jensen]{l'inégalité de
Jensen (p.~\pageref*{II-jensen})}, on a
\[\mathbb{E}_g\left(h^2(X)\frac{f^2(X)}{g^2(X)}\right) \geq \mathbb{E}_g\left(|h(X)|\frac{f(X)}{g(X)}\right)^2 = \left(\int|h(x)|f(x)dx\right)^2,\]
qui nous donne une borne inférieure indépendante du choix de \(g\). Elle
est atteinte en prenant \(g=g^\ast\).\hfill$\blacksquare$
\end{preuve}

Ce résultat formel n'a que peu d'intérêt pratique~: le choix optimal de
\(g\) fait intervenir \(\int|h(x)|f(x)dx\), qui est à une valeur absolue
près la quantité que l'on souhaite estimer\,! Il suggère néanmoins de
considérer des densités \(g\) telles que \(|h|f/g\) est quasi constante
et de variance finie. On se reportera au chapitre 3 de C. P. Robert and
G. Casella (\protect\hyperlink{refV-roca}{2004}) pour des exemples où un
bon choix de \(g\) permet des améliorations considérables par rapport à
des estimateurs de Monte-Carlo plus naïfs.

\chapter{Exercices}\label{V-exercices}

\section{Loi uniforme dans un
domaine}\label{V-loi-uniforme-dans-un-domaine}

Ecrire un algorithme pour simuler un point uniforme dans les trois
domaines représentés ci-dessous.

\begin{figure}
\centering
\fbox{[figure manquante : ex1-eps-converted-to.pdf]}
\caption{Domaine A}\label{V-fig-ex1-eps-converted-to.pdf}
\end{figure}

\begin{figure}
\centering
\fbox{[figure manquante : ex2-eps-converted-to.pdf]}
\caption{Domaine B}\label{V-fig-ex2-eps-converted-to.pdf}
\end{figure}

\begin{figure}
\centering
\fbox{[figure manquante : ex3-eps-converted-to.pdf]}
\caption{Domaine C}\label{V-fig-ex3-eps-converted-to.pdf}
\end{figure}

\section{Simulation selon la loi
géométrique}\label{V-simulation-selon-la-loi-guxe9omuxe9trique}

Cette loi correspond à la loi du nombre d'essais avant d'obtenir un
évènement de probabilité \(p\)
\[X\sim{\cal G} (p)~~\mathbb{P}(X=n)=p(1-p)^{n-1}~~n=1,2,\ldots\]

\runinhead{Question
1}\phantomsection\addcontentsline{toc}{subsubsection}{Question
1}\label{V-loigeom1}
Calculer l'espérance et la variance de \(X\)\solref{V-answer-loigeom1}

\runinhead{Question
2}\phantomsection\addcontentsline{toc}{subsubsection}{Question
2}\label{V-loigeom2}
Générer un échantillon de taille 1000 pour \(p=0.6\), en utilisant une
approche génétique. Calculer la moyenne et la variance. Comparer avec
les valeurs théoriques.\solref{V-answer-loigeom2}

\runinhead{Question
3}\phantomsection\addcontentsline{toc}{subsubsection}{Question
3}\label{V-loigeom3}
Montrer que si \(X\sim{\cal E}(\lambda)\), alors
\(\lceil X \rceil \sim{\cal G}(p)\). On précisera la valeur de
\(\lambda\).\solref{V-answer-loigeom3}

\runinhead{Question
4}\phantomsection\addcontentsline{toc}{subsubsection}{Question
4}\label{V-loigeom4}
Donner, implémenter et tester un algorithme pour simuler \({\cal G}(p)\)
directement\solref{V-answer-loigeom4}

\section{Simulation de la loi gaussienne par la méthode de
rejet}\label{V-simulation-de-la-loi-gaussienne-par-la-muxe9thode-de-rejet}

\(X\) est une variable aléatoire de loi gaussienne d'espérance \(m=0\)
et de variance \(\sigma^2=1\), \(\mathcal{N}(0,1)\). Sa densité est
\[f(x) = \frac{1}{\sqrt{2\pi}}\exp\left(-\frac{1}{2}x^2\right),~x\in \mathbb{R}\]

\runinhead{Question
1}\phantomsection\addcontentsline{toc}{subsubsection}{Question
1}\label{V-lgrej1}
Montrer que \(f(x) \leq C \exp \left(-|x|\right)\) où
\(C=\sqrt{\frac{e}{2\pi}}\)\solref{V-answer-lgrej1}

\runinhead{Question
2}\phantomsection\addcontentsline{toc}{subsubsection}{Question
2}\label{V-lgrej2}
On considère
\(g(x) = \frac{1}{2}\exp \left(-|x|\right),~x\in \mathbb{R}\). Montrer
que \(g\) est une densité de probabilité.\solref{V-answer-lgrej2}

\runinhead{Question
3}\phantomsection\addcontentsline{toc}{subsubsection}{Question
3}\label{V-lgrej3}
Proposer un algorithme pour simuler selon \(g\).\solref{V-answer-lgrej3}

\runinhead{Question
4}\phantomsection\addcontentsline{toc}{subsubsection}{Question
4}\label{V-lgrej4}
Implémenter l'algorithme de rejet pour simuler selon
\(f\).\solref{V-answer-lgrej4}

\section{Simulation selon la loi de Wigner}\label{V-wigner}

La loi de Wigner (ou du demi-cercle) est la loi de support \([-2, 2]\)
et de densité \(f(x) = \frac{1}{2\pi}\sqrt{4 - x^2}\). Simuler un grand
nombre de variables aléatoires de loi de Wigner, représenter
l'histogramme associé et le comparer à la
densité.\solref{V-answer-wigner}

\section{Loi des grands nombres et théorème central
limite}\label{V-loi-des-grands-nombres-et-thuxe9oruxe8me-central-limite}

\runinhead{Question
1}\phantomsection\addcontentsline{toc}{subsubsection}{Question
1}\label{V-lgntcl1}
Afin d'illustrer la \href{Probabilité\%20IV.pdf\%20\#lfgn}{Loi des
Grands Nombres}, visualiser la suite
\(S_n = \frac{X_1 + \ldots +X_n}{n}\) pour
\((X_i)_{i \in \mathbb{N}^\ast}\) une suite de variables aléatoires
indépendantes de loi uniforme sur {[}-1, 1{]}.

Indication~: pour \(x= [x_1 , \ldots, x_n]\) vecteur,
\texttt{np.cumsum(x)} est le vecteur
\([x_1 , x_1 + x_2 , x_1 + x_2 + x_3 , \ldots , x_1 + \ldots + x_n]\)
des sommes cumulées des coordonnées de \(x\).\solref{V-answer-lgntcl1}

\runinhead{Question
2}\phantomsection\addcontentsline{toc}{subsubsection}{Question
2}\label{V-lgntcl2}
Faire de même avec des variables aléatoires \(Y_i\) de loi de Cauchy
(c'est-à-dire de densité \(\frac{1}{\pi(1+x^2)}\)). La suite
\(\frac{Y_1 + \ldots +Y_n}{n}\) semble-t-elle converger\,?
Pourquoi\,?\solref{V-answer-lgntcl2}

\runinhead{Question
3}\phantomsection\addcontentsline{toc}{subsubsection}{Question
3}\label{V-lgntcl3}
Calculer la moyenne \(\mu\) et l'écart-type \(\sigma\) des \(X_i\) et
vérifier, afin d'illustrer le
\href{Probabilité\%20IV.pdf\%20\#TCL}{Théorème Central Limite}, que
\(\sqrt{n} (S_n - \mu)/\sigma\) converge en loi, lorsque
\(n \to \infty\), vers une variable de loi \(\mathcal{N}(0, 1)\)~: on se
fixera une grande valeur de n, on simulera un grand nombre de fois
\(\sqrt{n} (S_n - \mu)/\sigma\), on en tracera l'histogramme et on le
comparera à la densité \((2\pi)^-1/2 e^{-x^2}\) la loi
\(\mathcal{N}(0, 1)\).\solref{V-answer-lgntcl3}

\section{Simulation d'un mélange de
gaussiennes}\label{V-simulation-dun-muxe9lange-de-gaussiennes}

On souhaite simuler selon la loi de mélange introduite dans l'exercice
Mélange de loi du chapitre 3.

\runinhead{Question
1}\phantomsection\addcontentsline{toc}{subsubsection}{Question
1}\label{V-melloi1}
Soit \(K\) une v.a.r. de loi~:
\[\mathbb{P}(K=x) = \left\{\begin{array}{ll} 1/2 & \text{ si }x = 1 \\ 1/4 & \text{ si }x = 2 \\ 1/4 & \text{ si }x = 3 \\ 0 & \text{ sinon} \end{array} \right.\]
Implémenter la méthode d'inversion pour simuler
\(K\)\solref{V-answer-melloi1}

\runinhead{Question
2}\phantomsection\addcontentsline{toc}{subsubsection}{Question
2}\label{V-melloi2}
Soit \(X_1\), \(X_2\) et \(X_3\) des variables aléatoires de lois
respectives \(\mathcal{N}(0,1)\), \(\mathcal{N}(5,1/2)\) et
\(\mathcal{N}(8,4)\). Implémenter un algorithme de simulation de
\(X_K\). Comparer l'histogramme obtenu pour un échantillon de taille
1000 avec la densité de \(X_K\).\solref{V-answer-melloi2}

\section{Echantillonnage
d'importance}\label{V-echantillonnage-dimportance-2}

On cherche à évaluer l'espérance d'une variable aléatoire \(X\)
gaussienne centrée réduite (de densité \(f_X\) et de f.d.r. \(F_X\))
sachant qu'elle dépasse la valeur 3.

\runinhead{Question
1}\phantomsection\addcontentsline{toc}{subsubsection}{Question
1}\label{V-echimp1}
Exprimer la densité conditionnelle de
\(X|X>3\).\solref{V-answer-echimp1}

\runinhead{Question
2}\phantomsection\addcontentsline{toc}{subsubsection}{Question
2}\label{V-echimp2}
Calculer la quantité souhaitée à partir d'un échantillon de taille 1000
généré selon \(F_X\). Quelle est la valeur du taux de
rejet\,?\solref{V-answer-echimp2}

\runinhead{Question
3}\phantomsection\addcontentsline{toc}{subsubsection}{Question
3}\label{V-echimp3}
Implémenter l'algorithme d'échantillonnage d'importance en prenant comme
loi instrumentale la loi gaussienne d'espérance 3 et de variance 1.
Quelle est la proportion de poids nuls\,?

\begin{center}\rule{0.5\linewidth}{0.5pt}\end{center}

\solref{V-answer-echimp3}


\chapter{Annexe}\label{V-annexe}

\section{Preuve de la méthode
d'inversion}\label{V-preuve-de-la-muxe9thode-dinversion}

Pour pouvoir démontrer le \protect\hyperref[V-invgen]{théorème de la
méthode d'inversion (p.~\pageref*{V-invgen})}, il faut d'abord établir
un certain nombre de propriétés de la réciproque généralisée d'une
fonction de répartition. Elle peuvent être visualisées sur la figure
ci-dessous.

\begin{proposition}\phantomsection\label{V-proprecgen}
Soit \(F\) une fonction de répartition. Alors sa réciproque généralisée
\(F^-\) satisfait les propriétés suivantes.

\begin{enumerate}
\def\labelenumi{\arabic{enumi}.}
\item
  \(F^-\) est croissante.
\item
  \(\forall\, x \in \mathbb{R}\)~: \(F^- \circ F (x) \leq x\).
\item
  \(\forall\, u \in\, ]0,1[\)~: \(F \circ F^-(u) \geq u\) avec égalité
  si \(u \in F(\mathbb{R})\).
\item
  \(\forall\, (u,x) \in\, ]0,1[\, \times \mathbb{R}\)~:
  \(\left\{F(x) \geq u\right\} \Leftrightarrow \left\{x \geq F^-(u)\right\}\)
  et
  \(\left\{F(x) < u\right\} \Rightarrow \left\{x \leq F^-(u)\right\}\).
\end{enumerate}
\end{proposition}

\begin{preuve}
Pour tout \(u \in\, ]0,1[\) on note
\(F^{-1}\left([u,1]\right) := \left\{x \in \mathbb{R}: F(x) \geq u\right\}\)
l'image réciproque de \([u,1]\) par \(F\).

\begin{enumerate}
\def\labelenumi{\arabic{enumi}.}
\item
  Soit \((u,v) \in\, ]0,1[^2\). Si \(u < v\) alors
  \(F^{-1}\left([v,1]\right) \subset F^{-1}\left([u,1]\right)\) d'où
  \(F^-(u) = \inf F^{-1}\left([u,1]\right) \leq \inf F^{-1}\left([v,1]\right) = F^-(v)\).
  La fonction \(F^-\) est donc bien croissante.
\item
  Soit \(x \in \mathbb{R}\), alors
  \(F^- \circ F(x) = \inf \bigl\{ z \in \mathbb{R}: F(z) \geq F(x) \bigr\} = \inf F^{-1}\left(\left[F(x),1\right]\right)\).
  Comme \(F\) est croissante sur \(\mathbb{R}\) on a
  \([x,+\infty[\, \subseteq F^{-1}\left(\left[F(x),1\right]\right)\)
  donc
  \(F^- \circ F(x) = \inf F^{-1}\left(\left[F(x),1\right]\right) \leq \inf [x,+\infty[\, = x\).
\item
  Soit \(u \in\, ]0,1[\).
\end{enumerate}

\begin{itemize}
\tightlist
\item
  Puisque \(F^{-1}\left([u,1]\right)\) est nécessairement non vide, il
  existe une suite décroissante
  \((x_n)_{n\in\mathbb{N}} \subseteq F^{-1}\left([u,1]\right)\)
  convergeant vers \(\inf F^{-1}\left([u,1]\right) = F^-(u)\). La
  croissance de \(F\) implique que la suite
  \(\bigl(F(x_n)\bigr)_{n\in\mathbb{N}}\) est elle aussi décroissante,
  minorée par \(u\) car
  \((x_n)_{n\in\mathbb{N}} \subseteq F^{-1}\left([u,1]\right)\) donc
  convergente. Sa limite est de même supérieure ou égale à \(u\). Comme
  \(F\) est continue à droite, cette dernière n'est autre que
  \[\lim_{n \rightarrow +\infty} F(x_n) = F\left(\lim_{n\rightarrow+\infty} x_n\right) = F \circ F^-(u).\]
  Nous avons donc bien \(F\circ F^-(u) \geq u\).
\item
  Supposons maintenant que \(u \in F(\mathbb{R})\). Alors
  \(F^{-1}\left(\{u\}\right) := \left\{x \in \mathbb{R}: F(x) = u\right\} \neq \varnothing\).
  Il existe donc une suite décroissante
  \((x_n)_{n\in\mathbb{N}} \subseteq F^{-1}\left(\{u\}\right)\)
  convergeant vers \(\inf F^{-1}\left(\{u\}\right) = F^-(u)\) par
  croissance de \(F\). Comme \(F\) est continue à droite en tout point
  de \(\mathbb{R}\) et \(F(x_n) = u\) pour tout \(n \in \mathbb{N}\), on
  a bien
  \[u = \lim_{n \rightarrow +\infty} F(x_n) = F\left(\lim_{n\rightarrow+\infty} x_n\right) = F \circ F^-(u).\]
\end{itemize}

\begin{enumerate}
\def\labelenumi{\arabic{enumi}.}
\setcounter{enumi}{3}
\tightlist
\item
  Soit \((u,x) \in\, ]0,1[\, \times \mathbb{R}\).
\end{enumerate}

\begin{itemize}
\tightlist
\item
  \textbf{Equivalence.} Supposons \(F(x) \geq u\). On a
  \(F^- \circ F(x) \geq F^-(u)\) par croissance de \(F^-\) (propriété
  1.) et \(F^- \circ F(x) \leq x\) d'après la propriété 2.
  Réciproquement, supposons que \(x \geq F^-(u)\). Alors par croissance
  de \(F\) on a \(F(x) \geq F \circ F^-(u)\) puis
  \(F \circ F^-(u) \geq u\) d'après la propriété 3.
\item
  \textbf{Implication.} Supposons \(F(x) < u\). Tout
  \(z \in F^{-1}\left([u,1]\right)\) vérifie \(F(z) \geq u > F(x)\).
  Comme \(F\) est croissante sur \(\mathbb{R}\) on a donc \(z \geq x\),
  ce qui implique \(x \leq \inf F^{-1}\left([u,1]\right) = F^-(y)\).
\end{itemize}

\hfill$\blacksquare$
\end{preuve}

\begin{figure}
\centering
\includegraphics{RecGen.tex.pdf}
\caption{Réciproque généralisée d'une fonction de
répartition}\label{V-fig-RecGen.tex}
\end{figure}

Nous pouvons maintenant établir la preuve du
\protect\hyperref[V-invgen]{théorème de la méthode
d'inversion (p.~\pageref*{V-invgen})}.

\begin{preuve}
Soient \(x \in \mathbb{R}\) et \((\Omega,\mathcal{A},\mathbb{P})\)
l'espace probabilisé sur lequel sont définies \(U\) et \(X\). D'après la
propriété 4 ci-dessus, on a
\(\left\{\omega \in \Omega : F_X^-(U(\omega)) \leq x \right\} = \left\{\omega \in \Omega : U(\omega) \leq F(x) \right\}\),
d'où
\[\mathbb{P}\left( F^-_X(U) \leq x \right) = \mathbb{P}\left(U \leq F_X(x)\right) = F_X(x).\]

\hfill$\blacksquare$
\end{preuve}

\chapter{Références}\label{V-ruxe9fuxe9rences}

\hypertarget{refs}{}
\begin{cslreferences}
\leavevmode\hypertarget{refV-roca}{}%
C. P. Robert, and G. Casella. 2004. \emph{Monte-Carlo Statistical
Methods}. 2nd edition, Berlin: Springer.

\leavevmode\hypertarget{refV-BoxMuller}{}%
G.E.P. Box, and M.E. Muller. 1958. ``A Note on the Generation of Random
Normal Deviates.'' \emph{Ann. Math. Stat.} 29: 610--11.

\leavevmode\hypertarget{refV-Lehmer}{}%
Lehmer, Derrick H. 1951. ``Mathematical Methods in Large-Scale Computing
Units.'' \emph{Annu. Comput. Lab. Harvard Univ.} 26: 141--46.

\leavevmode\hypertarget{refV-Ziggurat}{}%
G. Marsaglia, and W. W. Tsang. 2000. ``The Ziggurat Method for
Generating Random Variables.'' \emph{Journal of Statistical Software}
5 (8): 1--7. \url{https://doi.org/10.18637/jss.v005.i08}.

\leavevmode\hypertarget{refV-ripley}{}%
Ripley, Brian D. 1987. \emph{Stochastic Simulation}. New York: Wiley.

\leavevmode\hypertarget{refV-vonNeumann}{}%
Von Neumann, John. 1951. ``Various Techniques Used in Connection with
Random Digits.'' \emph{U.S. Nat. Bur. Stand. Math. Ser.} 12: 36--38.
\end{cslreferences}
